% 第5章图：Householder反射的几何。解析式：y=(2,1)，||y||=sqrt5，按式(ls-householder)取 s=+1，
% v=y+||y|| e1=(2+sqrt5,1)≈(4.2361,1)，Hy=-||y|| e1≈(-2.2361,0)；y-Hy=v，中点(-0.1180,0.5)在 v^⊥ 上。
% 比例：1单位=12mm。超平面方向 (-1,2+sqrt5)/||·||≈(-0.2298,0.9733)。
\begin{tikzpicture}[x=1mm,y=1mm,>=Stealth,
  note/.style={font=\FigureNoteFont\fontsize{7.5}{9}\selectfont,text=FigureMuted}]
\def\u{12}
% e1 axis
\draw[FigureGrid,line width=.6pt,->] (-32,0)--(56,0) node[note,anchor=west] {$\bm e_1$};
% circle of radius ||y||
\draw[FigureGrid,line width=.6pt,densely dotted] (26.565:{2.2361*\u}) arc[start angle=26.565,end angle=180,radius={2.2361*\u}];
% hyperplane v^perp
\draw[FigureStroke,line width=.7pt,dashed] ({0.2298*1.45*\u},{-0.9733*1.45*\u})--({-0.2298*1.45*\u},{0.9733*1.45*\u})
  node[anchor=south,font=\figurefont\fontsize{8}{9.6}\selectfont] {$\bm v^\perp$};
\node[note,anchor=west] at ({0.2298*1.2*\u+1},{-0.9733*1.2*\u}) {反射超平面};
% reflection segment y -> Hy (parallel to v)
\draw[FigureMuted,line width=.6pt,densely dashed] ({2*\u},{1*\u})--({-2.2361*\u},0);
% right-angle mark at midpoint M=(-0.1180,0.5)
\begin{scope}[shift={({-0.1180*\u},{0.5*\u})}]
  \draw[FigureStroke,line width=.5pt] ({-0.2298*1.8},{0.9733*1.8})--++({0.9733*1.8},{0.2298*1.8})--++({0.2298*1.8},{-0.9733*1.8});
\end{scope}
% vectors
\draw[->,FigureGreen,line width=1.1pt] (0,0)--({4.2361*\u},{1*\u});
\draw[->,FigureBlue,line width=1.1pt] (0,0)--({2*\u},{1*\u});
\draw[->,FigureBlue,line width=1.1pt] (0,0)--({-2.2361*\u},0);
\fill[FigureStroke] (0,0) circle (.6);
% labels
\node[anchor=south,font=\figurefont\fontsize{8}{9.6}\selectfont] at ({2*\u},{1*\u+.6}) {$\bm y$};
\node[anchor=west,font=\figurefont\fontsize{8}{9.6}\selectfont,text=FigureGreen!70!black] at ({4.2361*\u+.8},{1*\u}) {$\bm v$};
\node[note,anchor=north] at ({3.2*\u},{0.755*\u-1.2}) {$\bm v=\bm y+\lVert\bm y\rVert_2\bm e_1$};
\node[anchor=north,font=\figurefont\fontsize{8}{9.6}\selectfont] at ({-2.2361*\u},-1) {$\bm H\bm y=-\lVert\bm y\rVert_2\bm e_1$};
\node[note,anchor=south east] at ({-1.25*\u},{0.32*\u+1}) {$\bm y-\bm H\bm y=\bm v$};
\end{tikzpicture}
